\documentclass[12pt,letterpaper]{article}
\usepackage{azzam}

\begin{document}

\section{Tuymaada Junior League 2025 Day 1}
\begin{enumerate}[resume]
    \item Positive integers $a$ and $b$ are given. A segment of positive integers contains more multiples of $a$ than multiples of $b$. Prove that it contains at least as many multiples of $2a$ as multiples of $2b$. 

    \item A quadrilateral $ABCD$ is given. Its diagonals are equal and intersect at a point $P$. The rays $AB$ and $DC$ intersect at $Q$. It is known that $\angle BPC = 2\angle BQC$. Prove that the circle passing through $B$ and tangent to the line $AC$ at $A$ is equal to the circle passing through $C$ and tangent to the line $BD$ at $D$. 

    \item In a zoo shop, four aviaries arranged in a circle contain 222 parrots each. Sometimes, a zoologist takes one parrot from an aviary and lets it go; together with it, she lets go either one parrot from the opposite aviary, or two parrots from the aviary on the left, or three parrots from the aviary on the right. At some moment, only one aviary still contains parrots. What is the least possible number of the remaining parrots? 

    \item There are $n > 3$ sportsmen in a tennis club. Last month every two of them played at most one match. It is known that if $A$ defeated $B$ and $B$ defeated $C$, then $C$ defeated $A$. Prove that it is possible to select at least $n/3$ sportsmen so that no two selected sportsmen played last month. 
\end{enumerate}

\section{Tuymaada Junior League 2025 Day 2}

\begin{enumerate}[resume]
    \item All positive real numbers are coloured with three colours (there is at least one number of each colour). Prove that there exist three numbers of three different colours so that each of them is less than the sum of the other two. 

    \item In a sequence $(x_n)$, the first number $x_1$ is positive, and
    \[ x_{n+1} = \frac{\{nx_n\}}{n} \quad \text{for} \quad n \geqslant 1 \]
    ($\{a\}$ denotes the fractional part of $a$). Prove that the sequence does not contain zeroes if and only if $x_1$ is irrational. 

    \item There are 128 persons in each of two rooms. In one move, we can select several (disjoint) pairs of persons so that the persons in each pair are in different rooms now, and exchange persons in each pair. What minimum number of moves is needed so that every two people found themselves in different rooms at least once? 

    \item All sides of a triangle $ABC$ are pairwise different. Its angle bisectors $AA_1$, $BB_1$, $CC_1$ meet at point $I$. The incircle of the triangle $ABC$ is tangent to the side $BC$ at $A_2$. The circle $\omega_a$ contains $A_1$, $A_2$, and the midpoint of $AI$. The circles $\omega_b$ and $\omega_c$ are defined similarly. Prove that the centres of $\omega_a$, $\omega_b$, $\omega_c$ are collinear. 
\end{enumerate}

\section{Tuymaada Junior League 2024 Day 1}
\begin{enumerate}
    \item \textit{Triangular numbers} are numbers of the form $1+2+\ldots+n$ with positive integral $n$, i.\,e. 1, 3, 6, 10, \dots . Determine the greatest non-triangular integer which is not a sum of different triangular numbers. 

    \item We call an \textit{edgehog} a graph in which one vertex is connected with all others and there are no other edges; the number of vertices of this graph we call the \textit{size} of the edgehog. A graph $G$ with $n > 1$ vertices is given. For each edge $e$, let us denote $s(e)$ the size of the largest edgehog in $G$ containing this edge. Prove the inequality (the summation is over all edges of $G$):
    \[ \sum_{e} \frac{1}{s(e)} \leqslant \frac{n}{2}. \]
    

    \item Three sportsmen ran with different constant velocities along a track of length 1. They started simultaneously at one end of the track. Upon reaching an end of the track, every sportsman immediately turned around and continued running in the opposite direction. After some time all three sportsmen met at the starting position and finished the practice. What is the maximum $S$ for which it is certainly true that at some moment the sum of pairwise distances between the sportsmen was at least $S$? 

    \item A triangle $ABC$ is given. $N$ and $M$ are the midpoints of $AB$ and $BC$, respectively. The bisector of angle $B$ meets the segment $MN$ at $E$. $H$ is the base of the altitude drawn from $B$ in the triangle $ABC$. The point $T$ on the circumcircle of $ABC$ is such that the circumcircles of $TMN$ and $ABC$ are tangent. Prove that points $T, H, E, B$ are concyclic. 
\end{enumerate}

\section{Tuymaada Junior League 2024 Day 2}

\begin{enumerate}[resume]
    \item In a $25 \times 25$ chessboard some squares are marked so that in each subboard of size $13 \times 13$ or $4 \times 4$ at least half the squares are marked. What is the minimum possible number of marked squares in the entire chessboard? 

    \item Extension of angle bisector $BL$ of the triangle $ABC$ (where $AB < BC$) meets its circumcircle at $N$. Let $M$ be the midpoint of $BL$. Isosceles triangle $BDC$ with base $BC$ and angle equal to $ABC$ at $D$ is constructed outside the triangle $ABC$. Prove that $CM \perp DN$. 

    \item Given are quadratic trinomials $f$ and $g$ with integral coefficients. For each positive integer $n$ there is an integer $k$ such that
    \[ \frac{f(k)}{g(k)} = \frac{n + 1}{n}. \]
    Prove that $f$ and $g$ have a common root. 

    \item A toy factory produces several kinds of clay toys. The toys are painted in $k$ colours. \textit{Diversity} of a colour is the number of \textit{different} toys of that colour. (Thus, if there are 5 blue cats, 7 blue mice and nothing else is blue, the diversity of colour blue is 2.) The painting protocol requires that \textit{each colour is used and the diversities of each two colours are different}. The toys in the store could be painted according to the protocol. However, a batch of clay Cheburashkas arrived at the store before painting (there were no Cheburashkas before). The number of Cheburashkas is not less that the number of the toys of any other kind. The total number of all toys, including Cheburashkas, is at least $\frac{(k+1)(k+2)}{2}$. Prove that now the toys can be painted in $k + 1$ colours according to the protocol. 
\end{enumerate}

\section{Tuymaada Junior League 2023 Day 1}

\begin{enumerate}[resume]
    \item[\textbf{1.}] Prove that for $a, b, c \in [0, 1]$ the following inequality holds:
    \[ (1 - a)(1 + ab)(1 + ac)(1 - abc) \leqslant (1 + a)(1 - ab)(1 - ac)(1 + abc). \]
    

    \item[\textbf{2.}] Serge and Tanya want to show Masha a magic trick. Serge leaves the room. Masha writes down a sequence $(a_1, a_2, \dots, a_n)$, where all $a_k$ equal 0 or 1. After that Tanya writes down a sequence $(b_1, b_2, \dots, b_n)$, where all $b_k$ also equal 0 or 1. Then Masha either does nothing or says ``Mutabor'' and replaces both sequences: her own sequence by $(a_n, a_{n-1}, \dots, a_1)$, and Tanya's sequence by $(1 - b_n, 1 - b_{n-1}, \dots, 1 - b_1)$. Masha's sequence is covered by a napkin, and Serge is invited to the room. Serge should look at Tanya's sequence and tell the sequence covered by the napkin. For what $n$ Serge and Tanya can prepare and show such a trick? Serge does not have to determine whether the word ``Mutabor'' has been pronounced. 

    \item[\textbf{3.}] Point $L$ inside triangle $ABC$ is such that $CL = AB$ and $\angle BAC + \angle BLC = 180^{\circ}$. Point $K$ on the side $AC$ is such that $KL \parallel BC$. Prove that $AB = BK$. 

    \item[\textbf{4.}] Two players play a game. They have $n > 2$ piles containing $n^{10} + 1$ stones each. A move consists of removing all the piles but one and dividing the remaining pile into $n$ nonempty piles. The player that cannot move loses. Who has a winning strategy, the player that moves first or his adversary? 
\end{enumerate}

\section{Tuymaada Junior League 2023 Day 2}

\begin{enumerate}[resume]
    \item[\textbf{5.}] A graph contains $p$ vertices numbered from 1 to $p$, and $q$ edges numbered from $p + 1$ to $p + q$. It turned out that for each edge the sum of the numbers of its ends and of the edge itself equals the same number $s$. It is also known that the numbers of edges starting in all vertices are equal. Prove that
    \[ s = \frac{1}{2}(4p + q + 3). \]
    

    \item[\textbf{6.}] An \textit{Euclidean step} transforms a pair $(a, b)$ of positive integers, $a > b$, to the pair $(b, r)$, where $r$ is the remainder when $a$ is divided by $b$. Let us call the \textit{complexity} of a pair $(a, b)$ the number of Euclidean steps needed to transform it to a pair of the form $(s, 0)$. Prove that if $ad - bc = 1$, then the complexities of $(a, b)$ and $(c, d)$ differ at most by 2. 

    \item[\textbf{7.}] $3n$ people forming $n$ families of a mother, a father and a child, stand in a circle. Every two neighbours can exchange places except the case when a parent exchanges places with his/her child (this is forbidden). For what $n$ is it possible to obtain every arrangement of those people by such exchanges? The arrangements differing by a circular shift are considered distinct. 

    \item[\textbf{8.}] Circle $\omega$ lies inside the circle $\Omega$ and touches it internally at point $P$. Point $S$ is taken on $\omega$ and the tangent to $\omega$ is drawn through it. This tangent meets $\Omega$ at points $A$ and $B$. Let $I$ be the centre of $\omega$. Find the locus of circumcentres of triangles $AIB$. 
\end{enumerate}


\end{document}
